\documentclass[12pt]{article}
\usepackage{amsmath,amssymb,amsfonts}
\usepackage[margin=1in]{geometry}
\usepackage{bm}

\newcommand{\vect}[1]{\mathbf{#1}}
\newcommand{\grad}{\boldsymbol{\nabla}}

\begin{document}

\section*{Problem 13}

\textbf{Problem:} The electromagnetic field tensor in free space. Problems 11 and 12 have laid the groundwork for the four-dimensional theory of the electromagnetic field. We shall use Gaussian units. To begin, if we identify the vector $H_\mu$ of Problem 11 with the magnetic field, then Problem 12 relates $H_\mu$ to a second-rank antisymmetric tensor $F_{\mu\nu}$.

\begin{itemize}
\item[(a)] Prove:
\[
F_{\mu\nu} = \partial_\mu A_\nu - \partial_\nu A_\mu, \qquad \text{where } \vect{A} = \vect{H} \times \vect{r}.
\]
Wait---more precisely, the problem asks us to define $F_{\mu\nu}$ in terms of the four-vector potential and show various properties. Let us work through the full problem as stated in the text.
\end{itemize}

\bigskip

The problem (as stated on pp.~40--41) has multiple parts. We present the full solution.

\subsection*{Setup}

In four-dimensional notation, we define the antisymmetric electromagnetic field tensor $F_{\mu\nu}$ ($\mu,\nu = 1,2,3,4$) such that:
\begin{itemize}
\item The spatial components ($i,j = 1,2,3$): $F_{ij} = \varepsilon_{ijk}H_k$ encodes the magnetic field $\vect{H}$.
\item The mixed components: $F_{i4} = -E_i$ and $F_{4i} = E_i$ encode the electric field.
\end{itemize}

Explicitly (where $A = \vect{H}$ and $\vect{E}$):
\[
F_{\mu\nu} = \begin{pmatrix}
0 & H_3 & -H_2 & -E_1 \\
-H_3 & 0 & H_1 & -E_2 \\
H_2 & -H_1 & 0 & -E_3 \\
E_1 & E_2 & E_3 & 0
\end{pmatrix}.
\]

\subsection*{(a) Show that $F_{\mu\nu}$ is antisymmetric}

By construction, $F_{\mu\nu} = -F_{\nu\mu}$:
\begin{itemize}
\item $F_{ij} = \varepsilon_{ijk}H_k = -\varepsilon_{jik}H_k = -F_{ji}$. \checkmark
\item $F_{i4} = -E_i$, $F_{4i} = E_i = -F_{i4}$. \checkmark
\item The four diagonal elements are zero. \checkmark
\end{itemize}

Only six of the remaining 12 elements are independent, expressing the fact that $\vect{E}$ and $\vect{H}$ have a total of six components.

\subsection*{(b) Writing $F_{\mu\nu}$ as a $4\times 4$ array in terms of $\vect{E}$ and $\vect{H}$}

As shown above:
\[
\boxed{F_{\mu\nu} = \begin{pmatrix}
0 & H_3 & -H_2 & -E_1 \\
-H_3 & 0 & H_1 & -E_2 \\
H_2 & -H_1 & 0 & -E_3 \\
E_1 & E_2 & E_3 & 0
\end{pmatrix}}.
\]

\subsection*{(c) Show directly from (a) that $F_{\mu\nu}$ is antisymmetric}

Done above. $F_{\mu\nu}$ has $4^2 = 16$ components. The four diagonal elements are zero. Only six of the remaining 12 elements are independent. So $F_{\mu\nu}$ has $\tfrac{1}{2}\cdot 4\cdot 3 = 6$ independent components, corresponding to the three components of $\vect{E}$ and three of $\vect{H}$.

We now extend to four dimensions.

\subsection*{(d) Making the identification $x_4 = ict$}

With the Minkowski four-vector $x_\mu = (x_1, x_2, x_3, x_4) = (x, y, z, ict)$ and the four-potential $A_\mu = (\vect{A}, i\varphi)$, we define:
\[
F_{\mu\nu} = \frac{\partial A_\nu}{\partial x_\mu} - \frac{\partial A_\mu}{\partial x_\nu} = \partial_\mu A_\nu - \partial_\nu A_\mu.
\]

The spatial components ($\mu = i$, $\nu = j$, with $i,j = 1,2,3$):
\[
F_{ij} = \partial_i A_j - \partial_j A_i = (\grad\times\vect{A})_k\,\varepsilon_{ijk} = \varepsilon_{ijk}B_k,
\]
where $\vect{B} = \grad\times\vect{A}$.

The mixed components ($\mu = i$, $\nu = 4$):
\[
F_{i4} = \partial_i A_4 - \partial_4 A_i = \partial_i(i\varphi) - \frac{\partial A_i}{\partial(ict)} = i\frac{\partial\varphi}{\partial x_i} - \frac{1}{ic}\frac{\partial A_i}{\partial t} = i\!\left(\frac{\partial\varphi}{\partial x_i} + \frac{1}{c}\frac{\partial A_i}{\partial t}\right) = -iE_i,
\]
since $\vect{E} = -\grad\varphi - \frac{1}{c}\frac{\partial\vect{A}}{\partial t}$.

This gives the full tensor structure consistent with the matrix above (with appropriate factors of $i$ from the Minkowski metric convention).

\subsection*{(e) Maxwell's equations in tensor form}

All four of Maxwell's equations can be expressed in two equations written in terms of the electromagnetic field tensor:

\textbf{Inhomogeneous equations} ($\grad\cdot\vect{D} = 4\pi\rho$ and $\grad\times\vect{H} = \frac{4\pi}{c}\vect{J} + \frac{1}{c}\frac{\partial\vect{D}}{\partial t}$):
\[
\boxed{\partial_\mu F_{\mu\nu} = \frac{4\pi}{c}\,J_\nu},
\]
where $J_\mu = (\vect{J}, ic\rho)$ is the four-current density.

\textbf{Homogeneous equations} ($\grad\cdot\vect{B} = 0$ and $\grad\times\vect{E} = -\frac{1}{c}\frac{\partial\vect{B}}{\partial t}$):
\[
\boxed{\partial_\lambda F_{\mu\nu} + \partial_\mu F_{\nu\lambda} + \partial_\nu F_{\lambda\mu} = 0},
\]
for all $\lambda,\mu,\nu$ (the Bianchi identity).

These equations are manifestly covariant---they take the same form in all inertial frames.

\subsection*{(f) Conservation of charge from Maxwell's equation}

From $\partial_\mu F_{\mu\nu} = \frac{4\pi}{c}J_\nu$, take $\partial_\nu$ of both sides:
\[
\partial_\nu\partial_\mu F_{\mu\nu} = \frac{4\pi}{c}\partial_\nu J_\nu.
\]

The left side vanishes because $\partial_\nu\partial_\mu$ is symmetric in $\mu\nu$ while $F_{\mu\nu}$ is antisymmetric. Therefore:
\[
\partial_\nu J_\nu = 0 \quad \Longleftrightarrow \quad \grad\cdot\vect{J} + \frac{\partial\rho}{\partial t} = 0.
\]

\subsection*{(g) The Lorentz force law}

The Lorentz force is:
\[
f_i = \frac{q}{c}\,F_{i\nu}\,\frac{dx_\nu}{ds},
\]
where $f_i$ are the three components of the force. Inserting the components of $F$:
\[
\vect{f} = q\!\left(\vect{E} + \frac{\vect{v}}{c}\times\vect{B}\right).
\]

The quantity $\frac{dx_\nu}{ds}$ is the four-velocity.

\subsection*{(h) Trace of the energy-momentum tensor}

The electromagnetic energy-momentum tensor is (see Problem 14):
\[
T_{\mu\nu} = \frac{1}{4\pi}\!\left(F_{\mu\lambda}F_{\nu\lambda} - \tfrac{1}{4}\delta_{\mu\nu}F_{\lambda\sigma}F_{\lambda\sigma}\right).
\]

Its trace is $T_{\mu\mu} = 0$ (shown in Problem 14). \hfill $\blacksquare$

\end{document}
